
seul polyn\^ome de Hilbert\,; donc (iii) r\'esulte de 2.11.

\textbf{Remarque 4.8.} Soit \(X\) propre et int\`egre sur \(S\) noeth\'erien. Sans hypoth\`ese de projectivit\'e sur \(X/S\), on ignore s'il reste vrai que \(\underline{\Pic}^\tau_{X/S}\to\underline{\Pic}_{X/S}\) est une immersion \uline{ferm\'ee}, sauf lorsque \(X\) est normal sur \(S\) et \(\underline{\Pic}_{X/S}\) est \uline{repr\'esentable} (FGA 236, 2.3).

% Added in pages 654--669
\providecommand{\NS}{\operatorname{NS}}
\providecommand{\Picard}{\operatorname{Pic}}
\providecommand{\Et}{\mathrm{et}}
\providecommand{\rk}{\operatorname{rk}}
\providecommand{\im}{\operatorname{Im}}
\providecommand{\HH}{\mathrm H}


\subsection*{5. Th\'eor\`emes de finitude du type ``N\'eron-S\'everi''.}

\textbf{Th\'eor\`eme 5.1.} \emph{Soit \(X\) propre sur \(S\) noeth\'erien\,; alors les groupes de N\'eron-S\'everi \(\underline{\Pic}_{X_K/K}(K)/\underline{\Pic}^0_{X_K/K}(K)\) des fibres g\'eom\'etriques \(X_K/K\) de \(X/S\) sont de type fini, et leur rang et l'ordre de leur sous-groupe de torsion \(\underline{\Pic}^\tau_{X_K/K}(K)/\underline{\Pic}^0_{X_K/K}(K)\) restent major\'es.}

En effet, par r\'ecurrence noeth\'erienne, on peut se borner \`a prouver 5.1 en y rempla\c{c}ant \(S\) par un ouvert non-vide arbitraire de \(S_{\red}\)\,; ainsi on peut supposer \(\underline{\Pic}^\tau_{X/S}\) repr\'esentable par un \(S\)-sch\'ema de type fini (XII 1.2 et 4.7 (iii)). Alors l'ordre de \(\underline{\Pic}^\tau_{X_K/K}(K)/\underline{\Pic}^0_{X_K/K}(K)\) est \'egal au nombre de composantes connexes de \((\Pic^\tau_{X/S})_K\), nombre qui reste le m\^eme au-dessus d'un ouvert non-vide de \(S\) (EGA IV 9.7.9).

Il reste \`a majorer le rang de \(\underline{\Pic}_{X_K/K}(K)/\underline{\Pic}^\tau_{X_K/K}(K)\). Comme dans la d\'emonstration de 4.7 (iii), on se ram\`ene facilement au cas \(S\) int\`egre, et \(X\) projectif et int\`egre sur \(S\). Soit \(r\) la dimension commune des fibres de \(X/S\). Si \(r\) est \(0\) (resp. \(1\)), alors ce rang est \'evidemment \(0\) (resp. \(1\)).

Si \(r\) est \(\ge2\), on va se ramener au cas \(r=2\) et \(X\) lisse sur \(S\). Quitte \`a restreindre \(S\) on peut trouver une section lin\'eaire \(Y\) de \(X\) qui est
